% ============================================================
% Sección 5: Conclusiones
% ============================================================

\section{Conclusiones}
\label{sec:conclusiones}

Este trabajo presentó una implementación de Redes Neuronales Informadas por Física
con activación sinusoidal (PINN-SIREN) para la simulación acelerada del speckle óptico
mediante un paradigma \textit{libre de malla} (\textit{mesh-free})
\citep{Karniadakis2021_review, Cuomo2022_review}.
Se resolvió la ecuación de Helmholtz 2D con campo complejo y condición de frontera de
fase aleatoria, eliminando la necesidad de discretización espacial del dominio.

Las conclusiones principales son:

\begin{enumerate}
  \item \textbf{Precisión excepcional en 1D y 2D.}
    La arquitectura SIREN alcanzó errores L2 de $0.006\%$ (1D) y $0.171\%$ (2D),
    ambos varios órdenes de magnitud por debajo del umbral de tesis ($< 5\%$).

  \item \textbf{La arquitectura es agnóstica a la condición de frontera.}
    La misma SIREN $5\times128$ resuelve Helmholtz tanto con onda plana suave (NB02)
    como con fase rugosa aleatoria (NB03), sin modificar la arquitectura ni la función
    de pérdida.
    Solo cambia el dato impuesto en $y = 0$.
    Esto demuestra la generalidad del método para cualquier perfil de rugosidad superficial.

  \item \textbf{Validación estadística del speckle.}
    El patrón generado satisface el criterio de \citet{Goodman2007_Speckle}
    para speckle completamente desarrollado: contraste $C = 1.0253$ ($|C-1| = 0.025 < 0.1$)
    y distribución de colas consistente con el valor teórico $e^{-2} \approx 0.135$.
    El test KS falla formalmente ($p < 0.0001$), atribuible a la alta potencia estadística
    con $N = 10{,}000$ puntos y a la ausencia de condición absorbente en $y = 1$.

  \item \textbf{La estrategia Adam + L-BFGS es esencial.}
    Adam proporciona exploración global robusta; L-BFGS refina la solución con
    información de curvatura de segundo orden.
    El uso exclusivo de cualquiera de los dos optimizadores produce resultados inferiores.

  \item \textbf{El peso físico $\lambda = 0.1$ equilibra datos y física en 2D.}
    En el caso 2D, la función de pérdida física incluye dos residuos (real e imaginario),
    duplicando efectivamente el peso de la física si se usa $\lambda = 1.0$.
    La calibración a $\lambda = 0.1$ fue crítica para la convergencia.
\end{enumerate}

\subsection*{Trabajo futuro}

Las extensiones naturales de este trabajo incluyen:

\begin{itemize}
  \item \textbf{Benchmark FEM} (NB04): comparación cuantitativa contra FEniCSx para
    múltiples realizaciones de speckle, cuantificando el \textit{Speed-up Factor}
    $S = T_{\mathrm{FEM}} / T_{\mathrm{PINN}}$ \citep{Thuerey2021_PBDL} y la
    huella de memoria pico (\textit{Peak Memory Footprint}), según los objetivos
    de la investigación.

  \item \textbf{Medios inhomogéneos} (NB05-NB06): extensión a la ecuación de Helmholtz
    con índice de refracción variable $n(\mathbf{r})$, relevante para speckle en
    materiales GRIN y tejidos biológicos.

  \item \textbf{Speckle en medios GRIN} (NB07): combinación de las extensiones anteriores
    para simular speckle en materiales con perfiles de índice gradiente,
    con potencial impacto en imágenes biomédicas y comunicaciones ópticas.

  \item \textbf{Condición de radiación de Sommerfeld}: incorporar
    $\mathcal{L}_{\mathrm{rad}} = |\partial E/\partial n - ikE|^2$
    en $y = 1$ para modelar propagación libre saliente, eliminando reflexiones espurias
    y mejorando la fidelidad estadística del speckle en el test KS.
\end{itemize}
